\section{Frontend: Interfaz Base y Autenticación de Cuentas}
\label{sec:s1_frontend}

El frontend materializa la capa de presentación de la identidad: layouts estructurados y
formularios interactivos de registro, login y gestión del perfil, soportados por el
\comando{authStore} de Zustand. El acceso a las vistas se filtra por la cadena de
\textit{guards} de la Figura~\ref{fig:s1_guards}.\par

\begin{figure}[!ht]
    \centering
    \begin{tikzpicture}[
        node distance=0.6cm,
        g/.style={draw=colorPrimario, fill=headerTabla, thick, rounded corners=2pt,
            text width=2.6cm, align=center, minimum height=1.0cm, font=\sffamily\scriptsize},
        d/.style={draw=grisISO, fill=grisTecnico, thick, rounded corners=2pt,
            text width=2.2cm, align=center, minimum height=0.9cm, font=\sffamily\scriptsize},
        rel/.style={-{Latex[length=2mm]}, colorPrimario, font=\sffamily\tiny}
    ]
        \node[g] (route) {Ruta solicitada};
        \node[g, right=of route] (prot) {ProtectedRoute\\(autenticación)};
        \node[g, right=of prot] (role) {RoleScopeGuard\\(rol)};
        \node[d, right=of role] (page) {Página por rol};
        \node[d, below=0.5cm of prot] (login) {Redirige a /login};

        \draw[rel] (route) -- (prot);
        \draw[rel] (prot) -- node[above]{ok} (role);
        \draw[rel] (role) -- node[above]{ok} (page);
        \draw[rel] (prot) -- node[right]{no auth} (login);
    \end{tikzpicture}
    \caption{Enrutamiento protegido del frontend: \texttt{ProtectedRoute} y \texttt{RoleScopeGuard}.}
    \label{fig:s1_guards}
\end{figure}

\subsection{Layouts: Navbar y Sidebar Responsiva}
\label{ssec:s1_fe_layouts}
Se implementó una \comando{Navbar} persistente y un componente \comando{Sidebar} responsivo
diseñado de forma exclusiva para renderizarse en entornos móviles y permanecer oculto en vistas
de escritorio. El enrutamiento protegido redirige según el rol tras el inicio de sesión.

\begin{itemize}[noitemsep]
    \item \textbf{Sesiones múltiples:} aislamiento por pestaña vía \comando{sessionStorage}.
    \item \textbf{Logout global:} propagación entre pestañas con \comando{BroadcastChannel}.
    \item \textbf{Redirección por rol:} PROFESSIONAL y RECRUITER acceden a paneles distintos.
    \item \textbf{Sidebar móvil:} se muestra bajo el \textit{breakpoint} móvil y se oculta en escritorio.
\end{itemize}

\subsection{Formularios Interactivos de Registro y Perfil}
\label{ssec:s1_fe_formularios}
Las páginas de \comando{pages/auth/} cubren login, registro, recuperación de contraseña y el
\textit{callback} OAuth2. Los formularios validan en cliente y muestran retroalimentación
inmediata; la entrada de OTP usa el componente \comando{input-otp}, y las notificaciones de
estado se presentan con \textit{toasts} (Sonner). La Figura~\ref{fig:s1_login_ui} muestra la
interfaz de autenticación.

\begin{figure}[!ht]
    \centering
    \fbox{\begin{minipage}[c][4.2cm][c]{0.6\textwidth}
        \centering
        \sffamily\color{grisISO}\footnotesize
        [\,Captura de pantalla pendiente\,]\par
        \vspace{0.3cm}
        Formulario de inicio de sesión / registro de EthosHub\par
        \vspace{0.2cm}
        \tiny(reemplazar por \texttt{assets/img/screenshots/login.png})
    \end{minipage}}
    \caption{Interfaz de autenticación de cuentas (vista de login/registro).}
    \label{fig:s1_login_ui}
\end{figure}

\begin{NotaTecnica}
El idioma del formulario se hidrata desde el perfil del usuario con \textit{fallback} estricto a
español, conforme a la estrategia i18n fijada en la Sección~\ref{sec:s0_frontend}. El
\comando{authStore} centraliza el estado de sesión y se limpia mediante \comando{resetAllStores}
en el cierre de sesión global.
\end{NotaTecnica}

\subsection{Stores y Servicios de Autenticación}
\label{ssec:s1_fe_stores}
La capa de presentación de la identidad se apoya en el \comando{authStore} y el
\comando{authService}, descritos en la tabla~\ref{tab:s1_fe_stores}.

\begin{table}[!ht]
    \centering
    \renewcommand{\arraystretch}{1.3}
    \fontsize{9}{10.8}\selectfont\sffamily
    \begin{tabularx}{\textwidth}{|l|X|}
        \hline
        \rowcolor{colorPrimario}
        \textcolor{white}{\textbf{Elemento}} & \textcolor{white}{\textbf{Responsabilidad}} \\ \hline
        \comando{authStore} & Sesión, usuario, rol y estado de autenticación. \\ \hline
        \rowcolor{grisTecnico}
        \comando{authService} & Llamadas a \comando{/api/auth/*} (registro, login, OAuth, recuperación). \\ \hline
        \comando{ProtectedRoute} & Bloquea rutas sin sesión y redirige a \comando{/login}. \\ \hline
        \rowcolor{grisTecnico}
        \comando{RoleScopeGuard} & Autoriza la vista según el rol del perfil. \\ \hline
        \comando{input-otp} & Componente de captura del código OTP. \\ \hline
        \rowcolor{grisTecnico}
        Sonner & \textit{Toasts} de retroalimentación de estado. \\ \hline
    \end{tabularx}
    \caption{Stores, servicios y componentes de autenticación (Sprint 1).}
    \label{tab:s1_fe_stores}
\end{table}

El aislamiento de sesión por pestaña (\comando{sessionStorage}) habilita varias cuentas
simultáneas; el \comando{BroadcastChannel} propaga el cierre de sesión a todas las pestañas,
invocando \comando{resetAllStores} para limpiar el estado global.

\clearpage
